\addtotoc{日本語要約}
\thispagestyle{plain}

% Tovchlol figures/tables are excluded from the global lists.
% textfont=up — япон фонтод налуу (italic) хэлбэр байхгүй тул тайлбарыг босоо болгоно
\captionsetup{list=no, textfont=up}

% Japanese caption labels for this section
\renewcommand{\figurename}{図}
\renewcommand{\tablename}{表}

\vspace*{-6mm}
\begin{center}
\textbf{\large 外国為替取引シグナル生成システム}\\[1mm]
\end{center}

\begin{flushright}
コンピュータサイエンス学科\\
学生：\studentID, M.~Munkhdorj\\
指導教員：N.~Soronzonbold
\end{flushright}

\vspace{1mm}
\setcounter{table}{0}
\setcounter{figure}{0}
\renewcommand{\thetable}{\arabic{table}}
\renewcommand{\thefigure}{\arabic{figure}}
\setlength{\columnsep}{7mm}
\begin{multicols*}{2}
\setlength{\parskip}{4pt}
\setlength{\intextsep}{4pt}
\setlength{\textfloatsep}{4pt}
\setlength{\abovecaptionskip}{2pt}
\setlength{\belowcaptionskip}{2pt}
\sloppy

%%─────────────────────────────────────────────
\noindent\textbf{1. 序論}

外国為替（FX）市場は変動が大きくノイズの多い時系列を生成するため、従来の統計的手法に基づく予測は限定的である。GBDT やニューラルネットワークのモデルはパターンを検出できることが示されているが、単一モデルのアーキテクチャは過学習のリスクにより実用面で依然として制約を受ける。本研究は、EUR/USD 通貨ペアに対して BUY/HOLD/SELL の分類シグナルを生成し、モバイルアプリを通じてユーザーへ配信する GBDT アンサンブル機械学習システムを構築する。3つの研究課題を設定した：
\begin{itemize}\setlength\itemsep{0pt}\small
    \item \textbf{RQ1：}モデルは時系列順を保持したホールドアウトテストで信頼性高く動作するか？
    \item \textbf{RQ2：}フィルタされたシグナルは実際の取引として約定するか？
    \item \textbf{RQ3：}固定リスクのルールの下で、1年間に正のリターンが得られるか？
\end{itemize}

%%─────────────────────────────────────────────
\noindent\textbf{2. 背景と関連研究}

金融時系列予測の研究は効率的市場仮説（EMH）の理論的懐疑のなかで発展してきたが、近年では機械学習が有意な程度にパターンを検出できることが裏付けられている。Krauss ら（2017）は GBDT が深層学習より信頼性の高い結果をもたらすと結論づけ、Ballings ら（2015）はアンサンブルが単一モデルより 3〜7\% 向上することを示した。Fischer・Krauss（2018）は LSTM を S\&P~500 で、Yildirim ら（2021）は EUR/USD で検証したものの、マルチタイムフレームの特徴量エンジニアリングを GBDT アンサンブルと組み合わせた実証研究は依然として少ない。

\textbf{貢献。}本研究は、GBDT アンサンブル、マルチタイムフレーム特徴量、信頼度しきい値によるシグナルフィルタリング、およびモバイルインターフェースを単一の統合システムとして FX 市場で実証的に検証した。

%%─────────────────────────────────────────────
\noindent\textbf{3. 研究手法}

モデル学習からモバイルアプリまでを貫く ML サービングシステムを構築した（図\ref{fig:tovjp_system_architecture}）。3つの中核レイヤーから成る：(i)~オフライン学習（GBDT pkl）；(ii)~サーバー（Flask + MongoDB + Redis + Expo Push）+ 4つの外部 API；(iii)~モバイルアプリ（React Native、4画面）。

\begin{figure}[H]
\centering
\adjustbox{max width=\columnwidth}{%
\begin{tikzpicture}[
    vm/.style={rectangle, rounded corners=5pt, draw=orange!85!black, fill=orange!22,
               line width=1.1pt, minimum width=4.8cm, minimum height=1.4cm,
               align=center, font=\normalsize\bfseries, inner sep=5pt},
    plainbox/.style={rectangle, rounded corners=4pt, draw=black!75, fill=white,
                     line width=1.0pt, minimum width=2.8cm, minimum height=1.0cm,
                     align=center, font=\normalsize, inner sep=4pt},
    websvc/.style={rectangle, rounded corners=5pt, draw=red!70!black, fill=red!25,
                   line width=1.0pt, minimum width=2.9cm, minimum height=1.0cm,
                   align=center, font=\normalsize\bfseries, text=black!85, inner sep=4pt},
    db/.style={cylinder, shape border rotate=90, aspect=0.30,
               draw=blue!65, fill=blue!18, line width=1.0pt,
               minimum width=2.4cm, minimum height=1.2cm,
               align=center, font=\normalsize\bfseries, text=blue!25!black},
    extapi/.style={rectangle, rounded corners=5pt, draw=green!55!black, fill=green!15,
                   line width=1.0pt, minimum width=3.0cm, minimum height=1.0cm,
                   align=center, font=\normalsize\bfseries, inner sep=4pt},
    pushsvc/.style={rectangle, rounded corners=5pt, draw=violet!70!black, fill=violet!18,
                    line width=1.0pt, minimum width=2.8cm, minimum height=1.1cm,
                    align=center, font=\normalsize\bfseries, inner sep=4pt},
    userbox/.style={rectangle, minimum width=0.9cm, minimum height=1.5cm,
                    draw=none, fill=none, inner sep=0pt},
    stickline/.style={line width=1.1pt, color=black!75, line cap=round},
    pipebox/.style={rectangle, rounded corners=4pt, draw=teal!60!black, fill=teal!14,
                    line width=1.0pt, minimum width=2.9cm, minimum height=1.0cm,
                    align=center, font=\normalsize, inner sep=4pt},
    flowarr/.style={-{Stealth[length=5pt 1.5, width'=0pt 0.6, sep=0pt 0.5]},
                    line width=1.0pt, color=black!72,
                    rounded corners=4pt, shorten >=2pt},
    thinarr/.style={-{Stealth[length=4pt 1.5, width'=0pt 0.6, sep=0pt 0.5]},
                    line width=0.7pt, color=black!60,
                    rounded corners=4pt, shorten >=2pt},
    busline/.style={line width=1.0pt, color=black!72},
    pusharr/.style={-{Stealth[length=5pt 1.5, width'=0pt 0.6, sep=0pt 0.5]},
                    line width=1.0pt, densely dashed, color=violet!65!black,
                    rounded corners=4pt, shorten >=2pt},
    deparr/.style={-{Triangle[open, length=4mm, width=3mm]},
                   line width=1.0pt, dashed, color=teal!55!black,
                   rounded corners=4pt, shorten >=2pt},
    elabel/.style={font=\small, text=black!65, fill=white, inner sep=1.5pt},
]
\node[userbox] (users) at (-9, 5.5) {};
\draw[stickline] ($(users.center)+(0, 0.46)$) circle (0.17);
\draw[stickline] ($(users.center)+(0, 0.29)$) -- ($(users.center)+(0, -0.25)$);
\draw[stickline] ($(users.center)+(-0.30, 0.08)$) -- ($(users.center)+(0.30, 0.08)$);
\draw[stickline] ($(users.center)+(0, -0.25)$) -- ($(users.center)+(-0.22, -0.65)$);
\draw[stickline] ($(users.center)+(0, -0.25)$) -- ($(users.center)+(0.22, -0.65)$);
\node[font=\small\bfseries, anchor=north] at ($(users.south)+(0, 0.05)$) {ユーザー};
\node[plainbox] (mob)  at (-4, 5.5) {モバイルアプリ\\[1pt]{\small (React Native)}};
\node[pushsvc]  (expo) at ( 0, 5.5) {Expo Push};
\draw[flowarr] (users.east) -- (mob.west);
\node[plainbox] (dns) at (-4, 3.5) {DNS};
\draw[flowarr] (mob.south) -- (dns.north);
\draw[black!55, line width=1.2pt, rounded corners=10pt]
    (-7, -5.6) rectangle (1, 2.4);
\node[font=\large\bfseries, anchor=north west, text=black!72]
    at (-6.85, 2.32) {サーバー};
\node[vm] (api)    at (-4,  0.9) {API サーバー\\[1pt]{\small Flask + Gunicorn}};
\node[vm] (worker) at (-4, -1.6) {ワーカー\\[1pt]{\small 推論・処理}};
\draw[flowarr] (dns.south)  -- (api.north);
\draw[flowarr] (api.south)  -- (worker.north);
\node[plainbox] (cache) at (-5.3, -4.2) {キャッシュ\\[1pt]{\small (Redis)}};
\node[db]       (mongo) at (-1.7, -4.4) {MongoDB};
\coordinate (db_trunk) at (-4, -2.7);
\draw[busline]  (worker.south) -- (db_trunk);
\draw[flowarr] (db_trunk) -| (cache.north);
\draw[flowarr] (db_trunk) -| (mongo.north);
\draw[thinarr] (cache.east) -- (-2.9, -4.2);
\fill[red!4, rounded corners=8pt]   (-12, -3.1) rectangle (-8, 1.95);
\draw[red!50, line width=0.8pt, rounded corners=8pt]
    (-12, -3.1) rectangle (-8, 1.95);
\node[font=\small\bfseries, anchor=north, text=red!55!black]
    at (-10, 1.85) {内部サービス};
\node[websvc] (mail)  at (-10,  0.9) {メール};
\node[websvc] (queue) at (-10, -1.1) {ジョブ\\キュー};
\node[websvc] (mon)   at (-10, -2.1) {監視};
\draw[thinarr] (mail.east) -- (api.west);
\draw[thinarr] (queue.east) -- ([yshift= 0.5cm]worker.west);
\draw[thinarr] (mon.east) -- ([yshift=-0.5cm]worker.west);
\fill[green!4, rounded corners=8pt]   (2.5, -2.9) rectangle (7.6, 3.2);
\draw[green!45!black, line width=0.8pt, rounded corners=8pt]
    (2.5, -2.9) rectangle (7.6, 3.2);
\node[font=\small\bfseries, anchor=north, text=green!40!black]
    at (5.05, 3.1) {サードパーティ API};
\node[extapi] (yf)  at (5.05,  1.9) {Yahoo Finance};
\node[extapi] (av)  at (5.05,  0.7) {AlphaVantage};
\node[extapi] (tv)  at (5.05, -0.5) {TradingView};
\node[extapi] (gem) at (5.05, -1.6) {Gemini API};
\draw[busline] (3.2, 1.9) -- (3.2, -1.6);
\draw[busline] (yf.west)  -- (3.2,  1.9);
\draw[busline] (av.west)  -- (3.2,  0.7);
\draw[busline] (tv.west)  -- (3.2, -0.5);
\draw[busline] (gem.west) -- (3.2, -1.6);
\node[circle, fill=black!72, minimum size=4pt, inner sep=0pt] at (3.2, -1.6) {};
\draw[flowarr] (3.2, -1.6) -- (worker.east);
\draw[pusharr] ([yshift=0.5cm]worker.east) -- ++(0.5, 0) -| (expo.south);
\draw[pusharr] (expo.west) -- (mob.east);
\fill[teal!4, rounded corners=8pt]   (-7, -8.7) rectangle (5, -6.5);
\draw[teal!50!black, line width=0.8pt, rounded corners=8pt]
    (-7, -8.7) rectangle (5, -6.5);
\node[font=\small\bfseries, anchor=north, text=teal!40!black]
    at (-1, -6.6) {オフライン学習};
\node[pipebox] (csv)   at (-5, -7.7) {履歴データ};
\node[pipebox] (train) at (-1, -7.7) {GBDT 学習};
\node[pipebox] (pkl)   at ( 3, -7.7) {model.pkl};
\draw[flowarr] (csv.east)   -- (train.west);
\draw[flowarr] (train.east) -- (pkl.west);
\node[elabel] at ([yshift=32pt, xshift=-20pt]pkl.north) {deploy};
\draw[deparr] (pkl.north) |- (1.5, -6.3) |- ([yshift=-0.5cm]worker.east);
\end{tikzpicture}%
}
\caption{システムアーキテクチャ}
\label{fig:tovjp_system_architecture}
\end{figure}

\textbf{データ。}EUR/USD の 2015〜2024 年の履歴データから、6つのタイムフレーム（M1〜H4）$\times$ 8つのテクニカル指標 = 48特徴量の行列を構築した。ラベルは将来240分の値動きに対する 30 pips のしきい値で BUY/SELL/HOLD に分類した。

\textbf{モデル。}LightGBM、XGBoost、CatBoost を等重みのソフト投票で統合したアンサンブル（図\ref{fig:tovjp_ensemble_arch}）。3つのモデルは成長戦略（leaf-wise、level-wise、symmetric）の違いにより互いを補完し、過学習のリスクを低減する。学習は WFV の第5フォールド（train 2015〜2022、val 2023、test 2024）で行った。過学習を防ぐため、木の深さ、低い学習率、早期終了、L1/L2 正則化、WFV、特徴量選択の6つの対策を組み合わせた。

\begin{figure}[H]
\centering
\adjustbox{max width=\columnwidth}{%
\begin{tikzpicture}[
    inp/.style ={rectangle, rounded corners=4pt, draw=blue!65, fill=blue!18,
                 line width=1.0pt, minimum width=6.0cm, minimum height=1.0cm,
                 align=center, font=\small\bfseries, text=blue!25!black},
    lgbm/.style={rectangle, rounded corners=4pt, draw=orange!85!black, fill=orange!22,
                 line width=1.0pt, minimum width=3.4cm, minimum height=1.25cm,
                 align=center, font=\small\bfseries, text=black!82},
    xgb/.style ={rectangle, rounded corners=4pt, draw=green!55!black, fill=green!15,
                 line width=1.0pt, minimum width=3.4cm, minimum height=1.25cm,
                 align=center, font=\small\bfseries, text=black!82},
    cat/.style ={rectangle, rounded corners=4pt, draw=teal!60!black, fill=teal!14,
                 line width=1.0pt, minimum width=3.4cm, minimum height=1.25cm,
                 align=center, font=\small\bfseries, text=black!82},
    vote/.style={rectangle, rounded corners=4pt, draw=violet!70!black, fill=violet!18,
                 line width=1.1pt, minimum width=6.0cm, minimum height=1.2cm,
                 align=center, font=\small\bfseries, text=black!82},
    outp/.style={rectangle, rounded corners=4pt, draw=red!70!black, fill=red!25,
                 line width=1.1pt, minimum width=3.4cm, minimum height=1.0cm,
                 align=center, font=\small\bfseries, text=black!85},
    sigb/.style={rectangle, rounded corners=4pt, draw=red!70!black, fill=red!25,
                 line width=1.1pt, minimum width=7.2cm, minimum height=1.0cm,
                 align=center, font=\small\bfseries, text=black!85},
    arr/.style ={-{Stealth[length=5pt 1.5, width'=0pt 0.6, sep=0pt 0.5]},
                 line width=1.0pt, color=black!72},
    lin/.style ={line width=1.0pt, color=black!72},
    lbl/.style ={font=\scriptsize\bfseries, fill=white, inner sep=1pt, text=black!65},
]
\node[inp]  (X)   at (0,  0)    {48 特徴量};
\node[lgbm] (lgb) at (-4, -2.4) {LightGBM};
\node[xgb]  (xgb) at ( 0, -2.4) {XGBoost};
\node[cat]  (cat) at ( 4, -2.4) {CatBoost};
\node[vote] (avg) at (0, -5.2)  {ソフト投票\\$\bar{p}_c = \tfrac{1}{3}\sum_i p_i^{(c)},\ c\in\{$BUY,HOLD,SELL$\}$};
\node[sigb] (sig)    at (0, -7.4)    {取引シグナル (BUY / HOLD / SELL、信頼度 $\in [0,1]$)};

\draw[arr] (X.south) -- (lgb.north);
\draw[arr] (X.south) -- (xgb.north);
\draw[arr] (X.south) -- (cat.north);

\coordinate (bus) at (0, -4.3);
\draw[lin] (lgb.south) -- ++(0, -0.4) -| (bus);
\draw[lin] (xgb.south) -- (bus);
\draw[lin] (cat.south) -- ++(0, -0.4) -| (bus);
\node[lbl, text=orange!85!black] at ($(lgb.south) + (0.35,-0.2)$) {$p_1$};
\node[lbl, text=green!55!black]  at ($(xgb.south) + (0.35,-0.2)$) {$p_2$};
\node[lbl, text=teal!60!black]   at ($(cat.south) + (0.35,-0.2)$) {$p_3$};
\draw[arr] (bus) -- (avg.north);

\draw[arr] (avg.south) -- (sig.north);
\end{tikzpicture}%
}
\caption{GBDT アンサンブルモデルのアーキテクチャ}
\label{fig:tovjp_ensemble_arch}
\end{figure}

\textbf{シグナルフィルタ。}信頼度（conf~$\geq 0.90$）と ATR（~$\geq 4$ pips）でシグナルをフィルタし、2025年の 359{,}639 件の予測から 3{,}567 件（1.0\%）が残り、MT5 の \texttt{MaxPositions=1} ルールにより 230 取引が約定した（図\ref{fig:tovjp_signal_funnel}）。SL/TP は ATR から動的に算出し、TP/SL=3:1 の比率を保持する。

\begin{figure}[H]
\centering
\adjustbox{max width=\columnwidth}{%
\begin{tikzpicture}[
    cnt/.style={font=\small\bfseries, text=black!82},
    pct/.style={font=\scriptsize, text=black!65},
    arr/.style={-{Stealth[length=5pt 1.5, width'=0pt 0.6, sep=0pt 0.5]},
                line width=1.0pt, color=black!72},
]
\node[trapezium, trapezium angle=78, trapezium stretches body,
      shape border rotate=180,
      draw=blue!65, fill=blue!18, line width=1.0pt,
      minimum height=1.0cm, minimum width=12cm,
      align=center, font=\small\bfseries, text=blue!25!black] (s1) at (0, 0)
    {1. 総予測数 (GBDT アンサンブル, 2025年)};
\node[cnt, right=0.5cm of s1] {359{,}639};
\node[pct, right=2.4cm of s1] {(100\%)};
\node[trapezium, trapezium angle=78, trapezium stretches body,
      shape border rotate=180,
      draw=orange!85!black, fill=orange!22, line width=1.0pt,
      minimum height=1.0cm, minimum width=8cm,
      align=center, font=\small\bfseries, text=black!82,
      below=0.5cm of s1] (s2)
    {2. 信頼度 $\geq 0.90$ フィルタ};
\node[cnt, right=0.5cm of s2] {$\sim$3{,}600};
\node[pct, right=2.2cm of s2] {($\approx$1.00\%)};
\node[trapezium, trapezium angle=78, trapezium stretches body,
      shape border rotate=180,
      draw=orange!85!black, fill=orange!22, line width=1.0pt,
      minimum height=1.0cm, minimum width=5.5cm,
      align=center, font=\small\bfseries, text=black!82,
      below=0.5cm of s2] (s3)
    {3. ATR $\geq 4.0$ pips フィルタ};
\node[cnt, right=0.5cm of s3] {3{,}567};
\node[pct, right=2.0cm of s3] {(0.99\%)};
\node[trapezium, trapezium angle=78, trapezium stretches body,
      shape border rotate=180,
      draw=teal!60!black, fill=teal!14, line width=1.0pt,
      minimum height=1.0cm, minimum width=3.6cm,
      align=center, font=\small\bfseries, text=black!82,
      below=0.5cm of s3] (s4)
    {4. MaxPositions=1\\{\scriptsize (MT5 EA 実行ルール)}};
\node[cnt, right=0.5cm of s4] {230};
\node[pct, right=1.6cm of s4] {(0.064\%)};
\node[rectangle, draw=red!70!black, fill=red!25, rounded corners=5pt,
      line width=1.1pt, minimum width=2.6cm, minimum height=0.9cm,
      align=center, font=\small\bfseries, text=black!85,
      below=0.5cm of s4] (s5)
    {5. 約定済み取引};
\node[cnt, right=0.5cm of s5] {230};
\node[pct, right=1.6cm of s5] {(TP: 103 / SL: 127)};
\draw[arr] (s1.south) -- (s2.north);
\draw[arr] (s2.south) -- (s3.north);
\draw[arr] (s3.south) -- (s4.north);
\draw[arr] (s4.south) -- (s5.north);
\node[pct, left=0.3cm of s2] {$-99.00\%$};
\node[pct, left=0.3cm of s3] {$-0.92\%$};
\node[pct, left=0.3cm of s4] {$-93.55\%$ (重複)};
\end{tikzpicture}%
}
\caption{シグナルフィルタリングのファネル（しきい値 0.90）}
\label{fig:tovjp_signal_funnel}
\end{figure}

\textbf{実験。}MT5 ストラテジーテスター、2025.01〜2026.01、初期資金 \$10{,}000、同時に1ポジション、実スプレッド + 10~point スリッページ、10個の信頼度しきい値（0.90〜0.99）。

%%─────────────────────────────────────────────
\noindent\textbf{4. 研究結果}

\textbf{分類性能。}未知の 2024 年テストでの全体正解率は \textbf{89.25\%}（全クラス）、検証（2023）サンプルでは conf~$\geq 0.90$ フィルタ後に正解率が \textbf{96.96\%} へ上昇した。検証−学習の差は +1.77\% で、過学習はみられない（表\ref{tab:tovjp_model_accuracy}）。

\begin{table}[H]
\centering
\caption{GBDT アンサンブルの正解率と F1 スコア}
\label{tab:tovjp_model_accuracy}
\scriptsize
\renewcommand{\arraystretch}{1.15}
\setlength{\tabcolsep}{3pt}
\begin{tabular}{|l|c|r|r|r|}
\hline
区分 & 期間 & Acc (\%) & F1-mac & F1-wgt \\ \hline
学習             & 2015〜2022 & 82.71 & 67.74 & 80.94 \\ \hline
検証             & 2023       & 84.48 & 62.67 & 82.29 \\ \hline
テスト (out-of-sample) & 2024  & \textbf{89.25} & \textbf{59.45} & \textbf{87.23} \\ \hline
\end{tabular}
\end{table}

\textbf{戦略リターン。}10個の信頼度しきい値（0.90〜0.99）すべてが正のリターンを示し、Buy~\&~Hold（+13.47\%）のベンチマークを 5.5〜38倍 上回った。0.91 のしきい値で最高の \textbf{+517.70\%（\$61{,}770）} に達し、0.90〜0.92 の帯域は高いリターンと低いドローダウン（MaxDD $<$~3\%）を両立するバランス点となった（図\ref{fig:tovjp_threshold_final_balance}）。

\begin{figure}[H]
\centering
\adjustbox{max width=\columnwidth}{%
\begin{tikzpicture}
\begin{axis}[
    width=0.95\textwidth,
    height=6.0cm,
    ybar,
    bar width=10pt,
    xlabel={信頼度しきい値},
    ylabel={最終残高 (\$)},
    xmin=0.895, xmax=0.995,
    ymin=0, ymax=70000,
    xtick={0.90,0.91,0.92,0.93,0.94,0.95,0.96,0.97,0.98,0.99},
    ytick={10000,20000,30000,40000,50000,60000},
    yticklabels={\$10K, \$20K, \$30K, \$40K, \$50K, \$60K},
    scaled y ticks=false,
    ymajorgrids=true,
    grid style={gray!30},
    tick label style={font=\scriptsize},
    label style={font=\small},
    legend style={at={(0.98,0.97)},anchor=north east,font=\scriptsize,draw=gray!40,fill=white},
    area legend,
    nodes near coords,
    nodes near coords style={font=\tiny,yshift=2pt,/pgf/number format/.cd,fixed,precision=0,1000 sep={,}}
]
\addplot[fill=blue!22, draw=blue!60!black] coordinates {
    (0.90,58262) (0.91,61770) (0.92,56512) (0.93,46807) (0.94,38050)
    (0.95,32844) (0.96,33040) (0.97,29019) (0.98,25447) (0.99,17429)
};
\end{axis}
\end{tikzpicture}%
}
\caption{しきい値ごとの最終残高}
\label{fig:tovjp_threshold_final_balance}
\end{figure}

\textbf{リスクと統計。}勝率 44.78〜65.71\%（Wilson 95\% CI）；二項検定 $z>5$、$p<0.0001$；$E[R]$=0.79〜1.63R ですべて正；MaxDD 1.98〜9.53\%（図\ref{fig:tovjp_risk_smallmult}）。Sharpe の 8.6〜12.8 という値は、サンプル数の少なさ、期間の短さ、コストの不完全なモデル化により\textbf{過大に算出されている}ため、しきい値間の相対比較にのみ有効である。

\begin{figure}[H]
\centering
\adjustbox{max width=\columnwidth}{%
\pgfplotsset{
    smallpanel/.style={
        width=7.2cm, height=4.4cm,
        xmin=0.895, xmax=0.995,
        xtick={0.90,0.92,0.94,0.96,0.98},
        tick label style={font=\tiny},
        label style={font=\scriptsize},
        title style={font=\scriptsize,yshift=-2pt},
        ymajorgrids=true, grid style={gray!30},
        ybar, bar width=4.5pt,
    }
}
\begin{tabular}{@{}cc@{}}
\begin{tikzpicture}
\begin{axis}[smallpanel, title={(A) MaxDD (\%)}, ylabel={MaxDD (\%)}, ymin=0, ymax=10.5]
\addplot[fill=red!28, draw=red!60!black] coordinates {
    (0.90,3.93) (0.91,2.97) (0.92,2.99) (0.93,9.53) (0.94,6.80)
    (0.95,3.94) (0.96,4.88) (0.97,2.96) (0.98,7.66) (0.99,1.98)};
\end{axis}
\end{tikzpicture} &
\begin{tikzpicture}
\begin{axis}[smallpanel, title={(B) Sharpe}, ylabel={Sharpe}, ymin=0, ymax=14]
\addplot[fill=blue!28, draw=blue!60!black] coordinates {
    (0.90,9.11) (0.91,9.75) (0.92,9.90) (0.93,9.21) (0.94,8.87)
    (0.95,8.57) (0.96,9.84) (0.97,9.02) (0.98,10.63) (0.99,12.83)};
\end{axis}
\end{tikzpicture} \\[4pt]
\begin{tikzpicture}
\begin{axis}[smallpanel, title={(C) Sortino}, ylabel={Sortino}, ymin=0, ymax=48]
\addplot[fill=teal!28, draw=teal!60!black] coordinates {
    (0.90,24.82) (0.91,25.99) (0.92,29.78) (0.93,25.92) (0.94,24.66)
    (0.95,23.97) (0.96,30.15) (0.97,24.90) (0.98,28.84) (0.99,44.76)};
\end{axis}
\end{tikzpicture} &
\begin{tikzpicture}
\begin{axis}[smallpanel, title={(D) Calmar = リターン / MaxDD}, ylabel={Calmar}, ymin=0, ymax=190]
\addplot[fill=orange!30, draw=orange!72!black] coordinates {
    (0.90,122.96) (0.91,174.20) (0.92,155.79) (0.93,38.63) (0.94,41.23)
    (0.95,57.97) (0.96,47.18) (0.97,64.19) (0.98,20.16) (0.99,37.53)};
\end{axis}
\end{tikzpicture}
\end{tabular}%
}
\caption{しきい値ごとのリスク指標}
\label{fig:tovjp_risk_smallmult}
\end{figure}

%%─────────────────────────────────────────────
\noindent\textbf{5. 結論}

本研究は、GBDT アンサンブル、信頼度しきい値フィルタリング、およびモバイルインターフェースを、学習からユーザーまでのエンドツーエンドのパイプラインとして EUR/USD 市場に統合し、実証的に検証した。3つの研究課題には次のように回答した：

\begin{itemize}\setlength\itemsep{0pt}\small
    \item \textbf{RQ1}：89.25\%（test 2024、全クラス）/ 96.96\%（val 2023、conf~$\geq 0.90$）
    \item \textbf{RQ2}：3{,}567 シグナル $\to$ 230 取引（6.45〜23.81\%）
    \item \textbf{RQ3}：+74〜518\% のリターン、$E[R]$=0.79〜1.63R
\end{itemize}

総合評価で上位の3つのしきい値（表\ref{tab:tovjp_rq3_summary}）：

\begin{table}[H]
\centering
\caption{総合評価上位3つのしきい値}
\label{tab:tovjp_rq3_summary}
\scriptsize
\renewcommand{\arraystretch}{1.12}
\setlength{\tabcolsep}{2.5pt}
\begin{tabular}{|c|r|r|r|l|}
\hline
しきい値 & リターン & MaxDD & 勝率 & プロファイル \\ \hline
0.91 & +517.70\% & 2.97\% & 47.17\% & \#1 バランス型 \\ \hline
0.92 & +465.12\% & 2.99\% & 47.26\% & \#2 バランス型 \\ \hline
0.99 & +74.29\%  & 1.98\% & 65.71\% & \#3 低リスク型 \\ \hline
\end{tabular}
\end{table}

これらの結果は、統計的な正解率にとどまらず、データ収集~$\to$~特徴量エンジニアリング~$\to$~アンサンブル学習~$\to$~シグナルフィルタリング~$\to$~自動約定~$\to$~モバイル通知という統合パイプラインの品質を併せて裏付けるものである。学術的には、GBDT アンサンブル、マルチタイムフレーム特徴量、信頼度しきい値フィルタリングを統合した実証的証拠を示し、実務的には、学習からモバイル配信までの完全な参照システムを構築し、リアルタイム環境で動作する能力を示した。

\textbf{制約。}結果は EUR/USD と 2025 年の期間でのみ検証されており、他の通貨ペアへ直接一般化することはできない；手数料（commission）とスワップ（swap）は完全にはモデル化されていない；シグナルから約定への変換率は 100\% 未満である。

\textbf{今後の課題。}GBP/USD、USD/JPY への拡張；LSTM/Transformer との比較；強化学習（RL）によるポジションサイズ管理；Model Registry と Google Play 配信の実装。

%%─────────────────────────────────────────────
\noindent\textbf{参考文献}

\begingroup
\scriptsize
\setlength{\parskip}{2pt}
\begin{enumerate}\setlength\itemsep{1pt}\setlength\parsep{0pt}\setlength\topsep{0pt}
    \item Krauss, C., Do, X.A., Huck, N. (2017). Deep Neural Networks, Gradient-Boosted Trees, Random Forests: Statistical Arbitrage on the S\&P~500. \emph{European Journal of Operational Research}, 259(2): 689--702.
    \item Ballings, M., Van den Poel, D., Hespeels, N., Gryp, R. (2015). Evaluating Multiple Classifiers for Stock Price Direction Prediction. \emph{Expert Systems with Applications}, 42(20): 7046--7056.
    \item Fischer, T., Krauss, C. (2018). Deep Learning with Long Short-Term Memory Networks for Financial Market Predictions. \emph{European Journal of Operational Research}, 270(2): 654--669.
    \item Ke, G., Meng, Q., Finley, T., Wang, T., Chen, W., Ma, W., Ye, Q., Liu, T.-Y. (2017). LightGBM: A Highly Efficient Gradient Boosting Decision Tree. \emph{Advances in Neural Information Processing Systems (NeurIPS)}, 30: 3146--3154.
    \item Yildirim, D.C., Toroslu, I.H., Fiore, U. (2021). Forecasting Directional Movement of Forex Data Using LSTM with Technical and Macroeconomic Indicators. \emph{Financial Innovation}, 7(1).
\end{enumerate}
\endgroup

\end{multicols*}
